\documentclass[11pt,a4paper]{article}
\usepackage{fontspec,amsmath,amssymb,amsthm,microtype}
\setmainfont{lmroman10-regular.otf}[BoldFont=lmroman10-bold.otf,ItalicFont=lmroman10-italic.otf,BoldItalicFont=lmroman10-bolditalic.otf]
\setmonofont{lmmono10-regular.otf}[Scale=MatchLowercase]
\usepackage[margin=25mm]{geometry}
\usepackage[unicode,hidelinks]{hyperref}
\setlength{\emergencystretch}{3em}
\setlength{\parskip}{0.4em}
\setlength{\parindent}{0pt}
\allowdisplaybreaks[2]
\raggedbottom
\title{\texorpdfstring{Couple ULA to an invariant Metropolis chain}{Couple ULA to an invariant Metropolis chain}}
\author{LangevinLean mathematical exposition}
\date{}
\begin{document}\maketitle
\noindent Source module: \texttt{Unadjusted\allowbreak{}Langevin/\allowbreak{}MALACoupling.\allowbreak{}lean}. The accompanying HTML pairs this mathematical exposition with the exact Lean statements and records the scope of the formalization.


A common Gaussian proposal couples unadjusted ULA with MALA. Under a stable small step, their expected distance grows by at most the Gibbs-average rejected displacement at each step. This is an auxiliary approximation for stationarity.


\section{\texorpdfstring{mala\allowbreak{}Acceptance}{malaAcceptance}}


Define \(a_h(x,y)=\min(1,e^{R_h(x,y)})\), with the proved target/\allowbreak{}proposal log ratio. The minimum caps acceptance at one and the exponential ensures nonnegativity; these bounds justify the accepted/\allowbreak{}rejected measure split.



{\raggedright\noindent Lean declarations: \texttt{mala\allowbreak{}Acceptance}.\par}

\section{\texorpdfstring{mala\allowbreak{}Proposal\allowbreak{}Acceptance}{malaProposalAcceptance}}


Compose \(a_h\) with \((x,\xi)\mapsto(x,T_h(x,\xi))\). The resulting \(a_h(x,T_h(x,\xi))\) is the acceptance weight on the original state/\allowbreak{}noise product space, prior to pushing to state/\allowbreak{}proposal pairs.



{\raggedright\noindent Lean declarations: \texttt{mala\allowbreak{}Proposal\allowbreak{}Acceptance}.\par}

\section{\texorpdfstring{mala\allowbreak{}Acceptance nonneg}{malaAcceptance nonneg}}


Both entries of the minimum, one and the exponential, are nonnegative. Their minimum is nonnegative. This direct order proof is used when combining real weights through \texttt{ENNReal.\allowbreak{}of\allowbreak{}Real}.



{\raggedright\noindent Lean declarations: \texttt{mala\allowbreak{}Acceptance\_\allowbreak{}nonneg}.\par}

\section{\texorpdfstring{mala\allowbreak{}Acceptance le one}{malaAcceptance le one}}


The minimum \(\min(1,e^{R_h})\) is at most its first entry. This proves \(a_h\le1\), so the rejection weight \(1-a_h\) is nonnegative.



{\raggedright\noindent Lean declarations: \texttt{mala\allowbreak{}Acceptance\_\allowbreak{}le\_\allowbreak{}one}.\par}

\section{\texorpdfstring{mala\allowbreak{}Proposal measurable}{malaProposal measurable}}


The gradient is continuous by Lipschitzness. Projections, scalar multiplication, subtraction, and addition preserve continuity, so the joint proposal map \((x,\xi)\mapsto x-h\nabla U(x)+\sqrt{2h}\xi\) is measurable.



{\raggedright\noindent Lean declarations: \texttt{mala\allowbreak{}Proposal\_\allowbreak{}measurable}.\par}

\section{\texorpdfstring{mala\allowbreak{}Proposal\allowbreak{}Acceptance measurable}{malaProposalAcceptance measurable}}


Differentiability makes \(U\) continuous, and Lipschitzness makes its gradient continuous. The log-ratio formula is assembled from continuous norms, inner products, and arithmetic; exponentiation and minimum preserve measurability. Composing with the measurable proposal gives the state/\allowbreak{}noise acceptance weight.



{\raggedright\noindent Lean declarations: \texttt{mala\allowbreak{}Proposal\allowbreak{}Acceptance\_\allowbreak{}measurable}.\par}

\section{\texorpdfstring{ula\allowbreak{}Law\allowbreak{}Step}{ulaLawStep}}


Define \(P_h\mu\) as the pushforward of \(\mu\otimes\gamma\) by the Euler proposal \(T_h\). The product measure encodes that the fresh Gaussian is independent of the current state. This transition is the actual independent-noise ULA step on laws.



{\raggedright\noindent Lean declarations: \texttt{ula\allowbreak{}Law\allowbreak{}Step}.\par}

\section{\texorpdfstring{mala\allowbreak{}One\allowbreak{}Step}{malaOneStep}}


Split \(\mu\otimes\gamma\) into acceptance- and rejection-weighted measures. Map the accepted part to \(T_h(x,\xi)\), and the rejected part to \(x\), then add them. This defines one full MALA transition with all rejection mass retained.



{\raggedright\noindent Lean declarations: \texttt{mala\allowbreak{}One\allowbreak{}Step}.\par}

\section{\texorpdfstring{mala\allowbreak{}Coupled\allowbreak{}Step}{malaCoupledStep}}


For a pair \((x,y)\), draw one shared Gaussian \(\xi\). ULA always moves to \(T_h(x,\xi)\); MALA moves to \(T_h(y,\xi)\) with weight \(a_h(y,T_h(y,\xi))\), otherwise stays at \(y\). The two weighted pushforwards define the explicit coupling measure.



{\raggedright\noindent Lean declarations: \texttt{mala\allowbreak{}Coupled\allowbreak{}Step}.\par}

\section{\texorpdfstring{weighted accept reject add}{weighted accept reject add}}


For a measurable \(a\in[0,1]\), combine the two with-density measures using \(a+(1-a)=1\). Nonnegativity of both terms justifies the \texttt{of\allowbreak{}Real} addition identity, so their sum is exactly the original measure.



{\raggedright\noindent Lean declarations: \texttt{weighted\_\allowbreak{}accept\_\allowbreak{}reject\_\allowbreak{}add}.\par}

\section{\texorpdfstring{mala\allowbreak{}Coupled\allowbreak{}Step fst}{malaCoupledStep fst}}


The first-coordinate update is identical on acceptance and rejection. Combine those two weights to one before pushing forward. The remaining product marginal is the original first marginal times a fresh Gaussian, so the first output law is \(P_h(\operatorname{fst}_*\rho)\).



{\raggedright\noindent Lean declarations: \texttt{mala\allowbreak{}Coupled\allowbreak{}Step\_\allowbreak{}fst}.\par}

\section{\texorpdfstring{mala\allowbreak{}Coupled\allowbreak{}Step snd}{malaCoupledStep snd}}


Project both branches of the explicit coupling to the second state. They are exactly the accepted proposal and rejected current-state branches of MALA. Product marginal identities replace \(\rho\) by its second marginal, giving \(\operatorname{snd}_*\rho'=M_h(\operatorname{snd}_*\rho)\).



{\raggedright\noindent Lean declarations: \texttt{mala\allowbreak{}Coupled\allowbreak{}Step\_\allowbreak{}snd}.\par}

\section{\texorpdfstring{mala\allowbreak{}Coupled\allowbreak{}Step is\allowbreak{}Probability}{malaCoupledStep isProbability}}


The first marginal is a measurable pushforward of a probability product by \texttt{mala\allowbreak{}Coupled\allowbreak{}Step\_\allowbreak{}fst}. It has mass one, and a measure and its first-coordinate pushforward have the same total mass. Therefore the joint coupling also has total mass one.



{\raggedright\noindent Lean declarations: \texttt{mala\allowbreak{}Coupled\allowbreak{}Step\_\allowbreak{}is\allowbreak{}Probability}.\par}

\section{\texorpdfstring{mala\allowbreak{}Coupled\allowbreak{}Chain}{malaCoupledChain}}


Initialize on the diagonal \((X_0,Y_0)=(X,X)\) with \(X\sim\mu\), then recursively apply the explicit coupled step. Equal starting states give zero initial transport cost, while subsequent first and second marginals follow ULA and MALA respectively.



{\raggedright\noindent Lean declarations: \texttt{mala\allowbreak{}Coupled\allowbreak{}Chain}.\par}

\section{\texorpdfstring{mala\allowbreak{}Coupled\allowbreak{}Chain is\allowbreak{}Probability}{malaCoupledChain isProbability}}


The diagonal pushforward of the initial probability is a probability. Induct on the number of transitions, using probability preservation of \texttt{mala\allowbreak{}Coupled\allowbreak{}Step} at each successor. This produces the instance needed for product noise and finite-cost arguments at every iteration.



{\raggedright\noindent Lean declarations: \texttt{mala\allowbreak{}Coupled\allowbreak{}Chain\_\allowbreak{}is\allowbreak{}Probability}.\par}

\section{\texorpdfstring{mala\allowbreak{}Coupled\allowbreak{}Chain fst}{malaCoupledChain fst}}


At zero, diagonal initialization has first marginal \(\mu\). In the successor step apply \texttt{mala\allowbreak{}Coupled\allowbreak{}Step\_\allowbreak{}fst} and the induction hypothesis. The resulting recursion is precisely the iteration \(P_h^n\mu\).



{\raggedright\noindent Lean declarations: \texttt{mala\allowbreak{}Coupled\allowbreak{}Chain\_\allowbreak{}fst}.\par}

\section{\texorpdfstring{mala\allowbreak{}Coupled\allowbreak{}Chain snd}{malaCoupledChain snd}}


At zero the second marginal is \(\pi\). The successor marginal theorem gives \(M_h\) applied to the current second marginal; substitute the induction hypothesis and the supplied one-step invariance \(M_h\pi=\pi\). Thus the second coordinate remains stationary through the coupling.



{\raggedright\noindent Lean declarations: \texttt{mala\allowbreak{}Coupled\allowbreak{}Chain\_\allowbreak{}snd}.\par}

\section{\texorpdfstring{mala\allowbreak{}Proposal nonexpansive}{malaProposal nonexpansive}}


Shared noise cancels in \(T_h(x,\xi)-T_h(y,\xi)\). The sharp strong-convexity/\allowbreak{}smoothness contraction theorem for the deterministic gradient map bounds the remaining norm by \(\|x-y\|\) when \(0\le h\le2/(m+L)\). No Gaussian moment is used in this pointwise comparison.



{\raggedright\noindent Lean declarations: \texttt{mala\allowbreak{}Proposal\_\allowbreak{}nonexpansive}.\par}

\section{\texorpdfstring{mala\allowbreak{}Coupled\allowbreak{}Step pointwise}{malaCoupledStep pointwise}}


The accepted distance is at most \(\|x-y\|\) by nonexpansiveness. On rejection, insert \(T_h(y,\xi)\) and use the triangle inequality to bound the distance by \(\|x-y\|+\|T_h(y,\xi)-y\|\). Weight the two branches; their common old-distance coefficients add to one, and the remaining term is the rejection displacement.



{\raggedright\noindent Lean declarations: \texttt{mala\allowbreak{}Coupled\allowbreak{}Step\_\allowbreak{}pointwise}.\par}

\section{\texorpdfstring{mala\allowbreak{}L1\allowbreak{}Cost}{malaL1Cost}}


Define \(\mathcal C(\rho)=\int\|x-y\|d\rho\) as a nonnegative extended integral. Keeping it extended avoids assuming integrability before the cost recurrence proves it. Subsequent theorems establish finiteness and then convert it to an ordinary expected distance.



{\raggedright\noindent Lean declarations: \texttt{mala\allowbreak{}L1\allowbreak{}Cost}.\par}

\section{\texorpdfstring{lintegral weighted map}{lintegral weighted map}}


Pull an extended nonnegative integral through a measurable pushforward, then through a with-density measure. Multiplication of the two nonnegative real factors becomes the embedded product \(a(x)b(f(x))\). This gives an exact integration formula for each branch of the coupling.



{\raggedright\noindent Lean declarations: \texttt{lintegral\_\allowbreak{}weighted\_\allowbreak{}map}.\par}

\section{\texorpdfstring{mala\allowbreak{}Coupled\allowbreak{}Step cost formula}{malaCoupledStep cost formula}}


Apply \texttt{lintegral\_\allowbreak{}weighted\_\allowbreak{}map} to accepted and rejected branches of the coupled measure. Their cost is the integral of the acceptance-weighted proposal distance plus rejection-weighted stay distance over \(\rho\otimes\gamma\). The formula exposes exactly the pointwise expression bounded by \texttt{mala\allowbreak{}Coupled\allowbreak{}Step\_\allowbreak{}pointwise}.



{\raggedright\noindent Lean declarations: \texttt{mala\allowbreak{}Coupled\allowbreak{}Step\_\allowbreak{}cost\_\allowbreak{}formula}.\par}

\section{\texorpdfstring{mala\allowbreak{}Coupled\allowbreak{}Step cost le}{malaCoupledStep cost le}}


Integrate the pointwise coupling bound. The old distance integrates to \(\mathcal C(\rho)\); the error depends only on the second state and innovation. Its second marginal is Gibbs by hypothesis, so that integral is bounded by \(C_{U,L}h\sqrt h\). Thus \(\mathcal C(\rho')\le\mathcal C(\rho)+C_{U,L}h\sqrt h\), valid initially as an extended inequality.



{\raggedright\noindent Lean declarations: \texttt{mala\allowbreak{}Coupled\allowbreak{}Step\_\allowbreak{}cost\_\allowbreak{}le}.\par}

\section{\texorpdfstring{mala\allowbreak{}Coupled\allowbreak{}Chain cost le}{malaCoupledChain cost le}}


Induct from zero diagonal cost. The second marginal remains Gibbs by the supplied one-step MALA invariance, so the one-step cost bound applies at every iteration. Summing the same local error gives \(\mathcal C(\rho_n)\le nC_{U,L}h\sqrt h\).



{\raggedright\noindent Lean declarations: \texttt{mala\allowbreak{}Coupled\allowbreak{}Chain\_\allowbreak{}cost\_\allowbreak{}le}.\par}

\section{\texorpdfstring{mala\allowbreak{}Coupled\allowbreak{}Chain distance integrable}{malaCoupledChain distance integrable}}


The quantitative extended cost estimate bounds \(\int\|x-y\|d\rho_n\) by a finite embedded real number. The distance is measurable and nonnegative, so finite extended integral is equivalent to ordinary integrability. This discharges the moment condition rather than assuming it.



{\raggedright\noindent Lean declarations: \texttt{mala\allowbreak{}Coupled\allowbreak{}Chain\_\allowbreak{}distance\_\allowbreak{}integrable}.\par}

\section{\texorpdfstring{mala\allowbreak{}Coupled\allowbreak{}Chain distance integral le}{malaCoupledChain distance integral le}}


Convert the finite extended cost inequality to real scalars using \texttt{ENNReal.\allowbreak{}to\allowbreak{}Real}. The newly proved integrability identifies the real value of the extended integral with the ordinary expected distance. Consequently \(\mathbb E_{\rho_n}\|X-Y\|\le nC_{U,L}h\sqrt h\).



{\raggedright\noindent Lean declarations: \texttt{mala\allowbreak{}Coupled\allowbreak{}Chain\_\allowbreak{}distance\_\allowbreak{}integral\_\allowbreak{}le}.\par}

\end{document}
